\section{Conclusion}
\label{sec:conclusion}

This volume has shown that gravitational dynamics can be interpreted as
arising from closure geometry.
Beginning from the modified distance definition
$w_{ab} = -\log S_{ab} + \epsilon L_{ab}$,
we have identified:

\begin{itemize}
  \item gravity as a modification of distance, not a new force
        (Section~\ref{sec:gravity_distance});
  \item the closure tensor $L_{\mu\nu}$ as the source of metric distortion
        (Section~\ref{sec:closure_density});
  \item the Einstein equation as the linear approximation of a fixed-point structure
        (Section~\ref{sec:einstein});
  \item the Schwarzschild solution as the exterior limit of closure geometry
        (Section~\ref{sec:schwarzschild});
  \item the black hole interior as a non-equilibrium closure phase with $\sigma > 0$
        (Section~\ref{sec:black_hole});
  \item cosmic expansion as the global $\sigma > 0$ regime,
        structurally analogous to the Friedmann equation
        (Section~\ref{sec:cosmology}).
\end{itemize}

The overall logical chain is:
\begin{equation}
  \text{Difference}
  \;\to\; C_{ij}
  \;\to\; J
  \;\to\; \omega
  \;\to\; L
  \;\to\; g_{\mu\nu}
  \;\to\; \text{Einstein equation}
  \;\to\; \text{Universe}.
\end{equation}

Taken together with Vol.~1 and Vol.~2,
this completes the structural derivation of time, space, matter, and gravity
from a single axiom: \emph{there is difference}.

The interpretation is structural and qualitative in nature,
and does not yet constitute a predictive theory.
Quantitative correspondence with observation remains a goal for future work.
\begin{equation}
  \boxed{\text{The world is not closed. It exists by never fully closing.}}
\end{equation}
